Frobenius's theorem: the Frobenius kernel is normal
← All problems

frobenius_kernel_isNormal

Submitter: Kim Morrison.

Notes: For a Frobenius group G acting transitively and faithfully on X with |X| ≥ 2, non-trivial point stabilisers, and the Frobenius condition (no non-identity element fixes more than one point), the set {1} ∪ {g | g fixes no point} is a normal subgroup. The only known proof uses Frobenius's induction-of-characters argument; no purely group-theoretic proof has been found in over a century.

Source: G. Frobenius, Über auflösbare Gruppen IV, Sitzungsber. Akad. Wiss. Berlin (1901).

Informal solution: Let H = stabilizer(x₀). Construct a class function θ on H of the form (1_H minus restriction of certain induced characters), and apply Frobenius reciprocity to lift θ to a generalised character of G whose kernel is exactly the Frobenius kernel K. The fact that the lift remains a virtual character (i.e. integer-valued combination of irreducibles) is exactly the content; that K is then a subgroup follows from K being a kernel.

theorem frobenius_kernel_isNormal (G X : Type) [Group G] [Fintype G] [Fintype X]
    [MulAction G X] [FaithfulSMul G X]
    (hcard : 2 ≤ Fintype.card X)
    (htrans : ∀ x y : X, ∃ g : G, g • x = y)
    (hstab : ∀ x : X, MulAction.stabilizer G x ≠ ⊥)
    (hfrob : ∀ g : G, g ≠ 1 → ∀ x y : X, g • x = x → g • y = y → x = y) :
    ∃ N : Subgroup G, N.Normal ∧
      (N : Set G) = {1} ∪ {g : G | ∀ x : X, g • x ≠ x} := by
  sorry